Queremos resolver o problema dual que foi proposta na \textbf{Pergunta 1}, isto é,
sejam $I^{-}: \mathbb{P} \rightarrow \mathbb{V}$ uma imagem antes de sofrer
alguma transformação, $\mathscr{T}^{-}_{\text{max}}$ a \textit{max tree} associada a $I^{-}$, 
$C \subset \mathbb{P}$ um conjunto conexo, $a = \text{\infimo{\le}}I^{-}(C)$ e 
$b \in \mathbb{V}$ tal que $b \le \text{\supremo{\le}}I^{-}(C)$.
A imagem $I^{+}$ é definida como:

$$\begin{array}{|llll}
   I^{+} : & \mathbb{P} & \rightarrow & \mathbb{V}  \\
           & x \in \mathbb{P} \setminus C & \rightarrow & I^{-}(x) \\
           & x \in C & \rightarrow & b
\end{array}$$

\begin{tcolorbox}[colback=gray!15, colframe=gray!50, arc=2mm]
  \textbf{Pergunta de Pesquisa:} Seja $\mathscr{T}_{\text{max}}^{+}$ a \textit{max tree} associada a $I^{+}$.
Sob essas condições, como devemos modificar $\mathscr{T}_{\text{max}}^{-}$
para obter $\mathscr{T}_{\text{max}}^{+}$?
\end{tcolorbox}

Resolver esse problema dual significa elaborar um algoritmo eficiente de
divisão de nós em uma \textit{max tree}. No entanto, para resolver esse
problema faz-se necessário responder perguntas de pesquisa auxiliares sobre 
modelagem teórica e topológica, complexidade algorítmica e estrutura de dados,
manutenção dinâmica de atributos, generalização e validação.

\begin{tcolorbox}[colback=gray!15, colframe=gray!50, arc=2mm]
\textbf{Pergunta sobre modelagem teórica:} Como definir formalmente a 
reestruturação das componentes conexas e a redistribuição de arestas 
quando a redução de valores de cinza em um conjunto de píxeis provoca a
 fragmentação de um nó da max-tree em múltiplos nós-filhos?
\end{tcolorbox}

\begin{tcolorbox}[colback=gray!15, colframe=gray!50, arc=2mm]
\textbf{Pergunta sobre complexidade algorítmica e estruturas de dados:}
Quais estruturas de dados e rotinas de rastreamento local permitem 
limitar a complexidade da operação de split, evitando que ela degrade 
para o custo de uma reconstrução $O(n\log n)$? 
\end{tcolorbox}

\begin{tcolorbox}[colback=gray!15, colframe=gray!50, arc=2mm]
\textbf{Pergunta sobre manutenção dinâmica de atributos:}
De que maneira os atributos dos nós (como área, momentos e caixas
delimitadoras) podem ser recalculados de forma local e incremental 
quando a massa de píxeis de um nó é particionada entre novas 
ramificações da árvore? 
\end{tcolorbox}

\begin{tcolorbox}[colback=gray!15, colframe=gray!50, arc=2mm]
\textbf{Pergunta generalização:}
Como combinar as operações de merge (já resolvidas) e de split em um 
modelo unificado que suporte alterações arbitrárias e simultâneas de 
níveis de cinza em pontos isolados da imagem? Existem outros operadores
necessários para generalizar? 
\end{tcolorbox}

\begin{tcolorbox}[colback=gray!15, colframe=gray!50, arc=2mm]
\textbf{Pergunta validação:}
Qual é o ganho de desempenho computacional obtido com o operador de 
split em filtros conectados complexos, granulometrias e no 
processamento interativo de imagens de grande escala em relação aos 
métodos tradicionais de reconstrução
\end{tcolorbox}